% ============================================================
%  Introduction.tex  –  Game Tracker
% ============================================================

\chapter{Introducción}

\section{Contexto}

En los últimos años, los videojuegos se han consolidado como uno de los medios
de entretenimiento más importantes a nivel mundial. Sin embargo, a medida que
las bibliotecas personales crecen, también lo hace la necesidad de organizarlas:
recordar qué títulos ya se completaron, cuáles están pendientes y cuál fue la
experiencia de juego resulta cada vez más difícil sin una herramienta de apoyo.

Las soluciones existentes suelen ser plataformas web de terceros que requieren
crear cuentas en servicios externos, o bien simples hojas de cálculo que no
ofrecen funcionalidad colaborativa. Game Tracker nace como respuesta a esta
necesidad: una aplicación móvil propia, ligera y enfocada exclusivamente en la
gestión de un catálogo de videojuegos con la posibilidad de escribir reseñas y
reaccionar a las opiniones de otros usuarios dentro de una misma comunidad.

\section{Descripción del proyecto}

Game Tracker es un sistema de gestión y reseñas de videojuegos compuesto por
dos componentes principales:

\begin{itemize}
    \item \textbf{Backend}: una API REST desarrollada en Java con Spring Boot,
    conectada a una base de datos MySQL. El servidor gestiona la lógica de
    negocio, la autenticación de usuarios con JWT y la persistencia de datos.
    Se desplegó en la plataforma Render, utilizando Aiven como servicio
    administrado de base de datos en la nube.

    \item \textbf{Frontend}: una aplicación móvil desarrollada en Flutter que
    consume la API REST. Permite a los usuarios explorar el catálogo, registrar
    nuevos juegos, escribir notas y reaccionar a las notas de otros usuarios
    desde su dispositivo móvil.
\end{itemize}

\section{Objetivos}

\subsection{Objetivo General}

Desarrollar una aplicación móvil funcional que permita a los usuarios gestionar
un catálogo personal de videojuegos, registrar su estado de progreso y compartir
reseñas dentro de una comunidad, aplicando patrones de diseño de software para
garantizar una arquitectura organizada y mantenible.

\subsection{Objetivos Específicos}

\begin{itemize}
    \item Implementar una API REST con Spring Boot que exponga operaciones
    CRUD sobre el catálogo de juegos y sus notas, protegida mediante
    autenticación con JWT y un sistema de tres roles.

    \item Desarrollar una interfaz móvil en Flutter que consuma la API y
    presente la información de forma clara e intuitiva para el usuario.

    \item Aplicar patrones de diseño reconocidos —MVC, Singleton, Strategy
    y ChangeNotifier— en la implementación del sistema para demostrar su
    utilidad práctica en un proyecto real.

    \item Desplegar el backend y la base de datos en servicios en la nube,
    garantizando que la aplicación sea accesible desde cualquier dispositivo
    móvil con conexión a internet.
\end{itemize}

\section{Alcance}

El proyecto abarca el ciclo completo de desarrollo de un sistema móvil con
backend en la nube: diseño de la base de datos, implementación del servidor,
desarrollo de la aplicación móvil y despliegue en producción. No se contempla
dentro del alcance la integración con plataformas externas de videojuegos como
Steam o PlayStation Network, ni la implementación de comunicación en tiempo real.